\section{An integral measure for all odd coefficients}
\label{sec:integral-beta-measure}

The dyadic beta multiplier encodes all the odd one-letter-depth
coefficients in a single integral measure. Its symmetry and distribution
identity will give their exact uniform divisibility.

Put $M=\ker\lambda$ and $N=\ker(\lambda,\rho_2^I)$.
For $g\in M$ write
\[
 B_g(X,Y)=1+(X-1)(Y-1)R_g(X,Y),\qquad
 q_g(X)=(X-1)R_g(X,1)=\sum_{n\ge1}a_n(g)(X-1)^n,
\]
where $a_n(g)=[S^nT]f_g$.  All coefficients lie in $\mathbf Z_2$.
The equality with the single-$T$ Magnus coefficients follows by taking
the metabelian image: the ending-$T$ part of $[x,y]$ is $S$, conjugation
by $x$ multiplies it by $X$, and a product of two commutator terms has
at least two occurrences of $T$.

The beta cocycle and its symmetry follow from Ihara's factorization
\cite[Proposition~1.6.1]{Ihara1999}, recalled in
\cite[Section~2.3, p.~101]{Furusho2003}. The boundary reflection is
an elementary consequence of the two cycle relations:
\[
 B(X,Y)B(XY,Z)=B(Y,Z)B(X,YZ),\quad B(X,Y)=B(Y,X),
 \quad q(X^{-1})=q(X).
\]
Here is a direct check of that reflection. The metabelian forms of
relations (I) and (II), at multiplier one, are
\[
 R(X,Y)=R(Y,X),\qquad
 R(X,Y)+(XY)^{-1}R(Y,(XY)^{-1})
       +Y^{-1}R((XY)^{-1},X)=0.
\]
The factors in the second formula are the commutator coefficients
of $[y,(xy)^{-1}]$ and $[(xy)^{-1},x]$, respectively.
Specializing at $Y=1$ and at $X=1$, and using symmetry, gives
$(X-1)R(X,X^{-1})=0$. Since the one-variable power-series ring is
a domain, $R(X,X^{-1})=0$ and $R(X^{-1},1)=-XR(X,1)$.
Multiplication by $X^{-1}-1$ proves $q(X^{-1})=q(X)$.

For $g\in N$, the normalized duplication equation is
\begin{equation}
 B_g(X,X)=B_g(-1,X).\label{eq:measure-duplication}
\end{equation}
The geometric input is \cite[Theorem~3.1 and the displayed
$(A_2^0)$, p.~102]{Furusho2003}: relations (I), (II), (IV) and
$(K_2)$ imply $B(X,X)=X^{-2\Psi_2}B(-1,X)$.
The convention there is $\Psi_2=-\rho_2^I$
\cite[p.~100, footnote]{Furusho2003}.
The full profinite $(K_2)$ relation follows from
\cite[Section~8, Theorem~8.1, equation~(50), author p.~30]{Enriquez2007}.
Relation (IV) on full $G$ is the input cited in Section~1.

\begin{lemma}[Integral boundary measure]
\label{lem:integral-boundary-measure}
For every $g\in N$ there is an integral $\mathbf Z_2$-valued measure
$\nu_{g,+}$ on $1+4\mathbf Z_2$, of total mass zero, such that
\begin{equation}
 \int t^{2k}\,d\nu_{g,+}(t)
  =(1-2^{2k})\tau_{2k+1}(g)\qquad(k\ge1),
 \label{eq:integral-boundary-moments}
\end{equation}
where the normalization is exactly
\[
 c_{2k+1}(g)=[u^{2k}]q_g(e^u),\qquad
 \tau_{2k+1}(g)=\frac{(2k)!}{2}c_{2k+1}(g).
\]
The same formula defines an integral continuous homomorphism
$\tau_{2k+1}:M\to\mathbf Z_2$ with conjugation weight
$\lambda_2^{2k+1}$. Its value modulo two is independent of $k$:
\[
 \tau_{2k+1}(g)\equiv\tau_3(g)=a_2(g)\pmod2
 \qquad(g\in M).
\]
\end{lemma}
\begin{proof}
The isomorphism
$\mathbf Z_2[[\mathbf Z_2]]\simeq\mathbf Z_2[[X-1]]$ gives an integral
measure $\mu_g$ with transform $q_g(X)=\int X^t\,d\mu_g(t)$.
One can verify the isomorphism directly: the ideals
$(2^r,X^{2^s}-1)$ and $(2^r,(X-1)^n)$ give cofinal topologies.
For fixed $r,n$, the coefficients below degree $n$ of $X^{2^s}-1$
are divisible by $2^r$ for large $s$; conversely,
$(X-1)^{2^s}$ is divisible by two modulo $X^{2^s}-1$.
The finite group-ring coefficients give the measure on residue classes.
Continuous functions are integrated by uniform locally constant
approximation, so pointwise divisibility passes to the integral.

Write $D_Y\log B=B^{-1}Y\partial B/\partial Y$; no analytic
logarithm is intended.  Differentiating the beta cocycle in $Z$ at
$Z=1$ gives
\[
 D_Y\log B(X,Y)=q(XY)-q(Y).
\]
Symmetry gives the analogous $X$ derivative.  Differentiating
\eqref{eq:measure-duplication}, with the chain rule on the diagonal,
therefore gives
\begin{equation}
 q(X)+q(-X)=2q(X^2).\label{eq:boundary-doubling-measure}
\end{equation}
Substitution $X\mapsto-X$ is legitimate: $X-1$ becomes
$-2-(X-1)$, and every resulting coefficient converges $2$-adically.

Boundary reflection makes $\mu_g$ invariant under $t\mapsto-t$.
The transforms $q(-X)$ and $q(X^2)$ correspond respectively to
$(-1)^t\mu_g$ and to the pushforward $[2]_*\mu_g$.
Equation~\eqref{eq:boundary-doubling-measure} and the torsion-freeness
of the measure module imply
\[
 \mu_g|_{2\mathbf Z_2}=[2]_*\mu_g.
\]
If $\nu_g=\mu_g|_{\mathbf Z_2^\times}$ and
$\nu_{g,+}=\mu_g|_{1+4\mathbf Z_2}$, it follows that
\[
 2\int t^{2k}\,d\nu_{g,+}
 =\int t^{2k}\,d\nu_g
 =(1-2^{2k})\int t^{2k}\,d\mu_g.
\]
The first equality uses negation symmetry, which exchanges the two
odd residue classes modulo four.  Expansion of $e^{tu}$ shows that
$\tau_{2k+1}=\tfrac12\int t^{2k}\,d\mu_g$, proving
\eqref{eq:integral-boundary-moments}; its factor $1-2^{2k}$ is a unit.
Also $\mu_g(\mathbf Z_2)=q_g(1)=0$, hence
$\nu_g(\mathbf Z_2^\times)=0$ and $\nu_{g,+}(1+4\mathbf Z_2)=0$.

The integral measure $\mu_g$ and its negation symmetry exist on all
of $M$, using only the integral boundary series and reflection.
Modulo two, the even residue classes contribute zero to
$\int t^{2k}\,d\mu_g$, while the two odd residue classes modulo four
have equal mass. The moment is therefore even, so its half is
integral. Moreover $t^{2k}-t^2$ is divisible by four for every
$t\in\mathbf Z_2$. Taking the integral and dividing by two proves
the asserted congruence; the identity $\tau_3=a_2$ follows from
$a_1=0$. Each moment is a finite rational linear combination of
boundary coefficients, hence is continuous.

On $M$ the beta multipliers multiply.  Their conjugation law is
$B_{hgh^{-1}}(X,Y)=B_g(X^{\lambda_2(h)},Y^{\lambda_2(h)})$,
as follows from the semilinear action on the metabelian commutator
module, with the multipliers of $h$ and $h^{-1}$ cancelling.
Thus $q$ is additive on $M$ and
$q_{hgh^{-1}}(X)=\lambda_2(h)q_g(X^{\lambda_2(h)})$.
The coefficient of $u^{2k}$ has the asserted weight.
\end{proof}

\begin{lemma}[Uniform odd-coordinate upper bound]
\label{lem:uniform-odd-coordinate-bound}
If $g\in N$ and $a_2(g)=a_4(g)=\cdots=a_{2r-2}(g)=0$, then
\[
 a_{2r}(g)\in2^{2r+1}\mathbf Z_2.
\]
\end{lemma}
\begin{proof}
Reflection implies that $q(1+S)$ is a power series in
$S^2/(1+S)$.  This follows recursively by cancelling the first
nonzero even-degree term of a reflection-invariant series; an odd
first degree would have its coefficient equal to its negative.
Consequently the assumed lower even coefficients make $q$ start in
$S^{2r}$.  The measure in Lemma~\ref{lem:integral-boundary-measure}
has zero moments in degrees $0,2,\ldots,2r-2$.
For $t=2z+1$ put
\[
 P_r(t)=\prod_{j=0}^{r-1}\bigl(t^2-(2j+1)^2\bigr)
 =2^{2r}(2r)!\binom{z+r}{2r}.
\]
The binomial polynomial is $\mathbf Z_2$-valued, by continuity from
integers.  Hence
\[
 (1-2^{2r})\tau_{2r+1}
 =\int t^{2r}\,d\nu_{g,+}
 =\int P_r(t)\,d\nu_{g,+}
 \in2^{2r}(2r)!\mathbf Z_2.
\]
As $\tau_{2r+1}=\frac{(2r)!}{2}a_{2r}$, cancellation in
$\mathbf Q_2$ gives the asserted divisibility.
\end{proof}

